\section{Conclusion}
\label{sec:conclusion}

This paper asked three questions of a single fixed panel of 218{,}934 thirty-minute S\&P~500 realized-variance bars, against a single fixed benchmark---per-bar OLS on HAR and calendar features, QLIKE $0.13415$.

\emph{Where does the information live?} Everywhere and nowhere in particular. Every exogenous bucket improves on HAR with overwhelming statistical significance, yet the buckets overlap so heavily that the sum of individual gains is roughly twice the joint gain, and the model confidence set over feature sets collapses to the singleton joint bucket. Within the least-squares class, 96\% of explainable variance belongs to the six HAR features; the 359-column exogenous block adds a unique contribution smaller than that of a handful of mechanism-targeted engineered features on a per-feature basis. The individual-feature landscape repeats the pattern one level down: one dominant cheap feature, one partially redundant companion, a long tail of tiny effects, and a genuine null.

\emph{Does extraction require nonlinearity?} Yes, but boundedly. Tree ensembles beat the best penalized linear model by $0.00184$ QLIKE on the full feature set, and carry a pure-nonlinearity premium of $0.00297$ on HAR features alone. The premium is real, systematic across all nine feature sets, and not a tuning artifact---indeed, tuning is a non-lever from every angle we measured: causally tuned hyperparameters trail fixed defaults, and even a test-set oracle over penalties moves the incumbent by at most $0.00015$. What the trees exploit saturates at additive-plus-pairwise structure anchored to the session clock; deeper capacity, deep sequence models, and mixture-of-experts routing all fail to improve on it. The resulting ordering---information $>$ model class $>$ hyperparameters---is the paper's central practical law.

\emph{What algorithm should do the extracting?} With features frozen, exactness and anchoring matter more than expressiveness. Rank-1 rolling least squares makes the every-bar OLS benchmark computable at scale; the online elastic-net homotopy extends the same exactness to $\ell_1$ estimation; the causal tuning protocol converts ``tuning does not help'' from anecdote into a defensible causal claim. The spectral $k$-nearest-neighbors proposal is instructive precisely because its ablations are decisive: the OLS anchor is load-bearing, large-neighborhood averaging is load-bearing, and the manifold embedding---the novel part---hurts on every matched cell, with a measured out-of-sample extension error that explains why. What survives is an anchored analog search that beats the incumbent (QLIKE $0.13096$, DM $-17.6$) but trails every causally tuned parametric arm. In the dense-but-weak regime, nonparametric flexibility is a cost to be justified, not a default.

\paragraph{Open items.} We flag the incomplete edges of the evidence explicitly. (i) The bucket-attribution battery of Section~\ref{sec:marginal} uses a negligible-penalty surrogate for OLS with exogenous features; the exact $\alpha=0$ battery is a pending robustness run on existing machinery. (ii) The linear-versus-nonlinear gap of Section~\ref{sec:linvsnonlin} is measured against the shared incumbent only; pairwise Diebold--Mariano tests and a model confidence set over the 45 arms are pending. (iii) The winning spectral-$k$NN exogenous arm survives on one bucket while sibling buckets failed as singular; the failure needs diagnosis before the bucket ranking can be trusted, and the correction-weight axis is unexplored. (iv) The attribution results point outward: the highest-value next investment is not a larger model but mechanism-targeted data---auction-imbalance, order-flow, and dealer-positioning series aligned to the session-transition regime where the nonlinear signal concentrates.
